\begin{proofbox}
\textbf{\textcolor{CProof}{证明目标}}

证明以下两件事：
\[
 \boxed{ID=IE=IF=r},\qquad
 \boxed{S=\frac12r(a+b+c)=rp}.
\]
其中 $D,E,F$ 依次是 $I$ 向 $BC,CA,AB$ 作垂线的垂足，
$a=BC,b=CA,c=AB$，$p=(a+b+c)/2$。

\medskip
\par\noindent\begin{minipage}{\linewidth}
\textbf{\textcolor{CProof}{方法一：门槛最低（初中全等三角形 + 面积拼接）}}

\PAngleFigure
\end{minipage}\par
\textbf{核心想法：}先不用记“角平分线上的点到两边距离相等”，
而是\textbf{亲手证明这个等距性质}，再把三块面积相加。
\begin{enumerate}[leftmargin=*,itemsep=0.7em]
 \item \textbf{证明 $IE=IF$。}
 因为 $I$ 在 $\angle BAC$ 的角平分线上，所以
 $\angle EAI=\angle IAF$。由作图，
 $\angle AEI=\angle AFI=90^\circ$，两直角三角形 $\triangle AEI$ 和
 $\triangle AFI$ 的斜边都是 $AI$。因此它们有两个对应角相等、
 一条对应边相等，可判定全等，于是对应边
 \[
  IE=IF.
 \]
 \textbf{为什么做这一步？}先确定内心到 $CA$、$AB$ 的\textbf{垂直距离}相等。

 \item \textbf{证明 $IF=ID$。}
 因为 $I$ 也在 $\angle ABC$ 的角平分线上，
 同理，$\triangle BFI$ 与 $\triangle BDI$ 都是直角三角形，
 斜边同为 $BI$，且 $\angle FBI=\angle IBD$，所以它们全等，
 \[
  IF=ID.
 \]
 连同上一步，得到 $ID=IE=IF$。把这个共同长度记作 $r$。

 \item \textbf{理解为什么这就是内切圆半径。}
 以 $I$ 为圆心、$r$ 为半径画圆，
 它到三条边的\textbf{垂直距离都恰好是 $r$}，
 因而与三边分别在 $D,E,F$ 处相切；
 因 $I$ 在三角形内部，这就是内切圆。所以共同距离为内切圆半径。

 \item \textbf{把大三角形分成三块。}
\PIncenterFigure{2}{4}{6}{0.95}{a}{b}{c}{1}{连三个顶点分成三块；三条实线垂线分别给出对应的高。}

 连接 $IA,IB,IC$。这三条线段把原三角形无重叠地分为
 $\triangle IBC,\triangle ICA,\triangle IAB$。
 它们分别以 $BC=a,CA=b,AB=c$ 为底，
 对应的高是 $ID=IE=IF=r$，于是
 \[
   S_{\triangle IBC}=\frac12 ar,\quad
   S_{\triangle ICA}=\frac12 br,\quad
   S_{\triangle IAB}=\frac12 cr.
 \]
 \textbf{为什么可以用这些高？}三角形的高是顶点到\textbf{底边所在直线的垂线段}，不是 $IA,IB,IC$。

 \item \textbf{将三块面积相加。}
 \[
 \begin{aligned}
 S&=S_{\triangle IBC}+S_{\triangle ICA}+S_{\triangle IAB}\\
  &=\frac12 ar+\frac12 br+\frac12 cr\\
  &=\frac12 r(a+b+c)=rp.
 \end{aligned}
 \]
 由于 $a+b+c>0$，两边同除以 $(a+b+c)/2$ 得
 \[
  \boxed{r=\frac{2S}{a+b+c}=\frac Sp}.
 \]
\end{enumerate}
\textbf{门槛：}仅使用角平分线、直角三角形全等、面积公式、等式变形，初中生可以逐步理解。

\medskip
\par\noindent\begin{minipage}{\linewidth}
\textbf{\textcolor{CProof}{方法二：步骤最少（考场直接写的面积法）}}

\PIncenterFigure{2}{4}{6}{1}{a}{b}{c}{1}{面积直接相加：$S=ar/2+br/2+cr/2$。}
\end{minipage}\par
\textbf{前提：}已经学过“内心到三边的距离都是内切圆半径”。
\begin{enumerate}[leftmargin=*,itemsep=0.55em]
 \item 连接 $IA,IB,IC$；因为 $I$ 是内心，所以
 $ID=IE=IF=r$。
 \item 三块面积相加，\textbf{一行完成核心计算}：
 \[
  S=\frac12ar+\frac12br+\frac12cr
   =\frac12r(a+b+c).
 \]
 \item 等式两边除以 $(a+b+c)/2$，得到
 \[
   \boxed{r=\frac{2S}{a+b+c}}.
 \]
\end{enumerate}
\textbf{考场价值：}如果题目只求 $r$ 或 $S$，这是建议优先采用的写法；
不用计算 $I$ 的具体位置。

\medskip
\par\noindent\begin{minipage}{\linewidth}
\textbf{\textcolor{CProof}{方法三：常规证明（角平分线性质 + 垂足 + 分别求面积）}}

\PIncenterFigure{2}{4}{6}{1}{a}{b}{c}{2}{观察角平分线 $AI$ 两旁的直角三角形：$IE=IF$；\\[3pt]在顶点 $B$ 处同理得到 $IF=ID$。}
\end{minipage}\par
\textbf{核心想法：}先规范地推出各条高相等，再按三角形面积公式逐项计算。
\begin{enumerate}[leftmargin=*,itemsep=0.65em]
 \item \textbf{作垂线。}过内心 $I$ 作
 $ID\perp BC$、$IE\perp CA$、$IF\perp AB$。
 由角平分线上的点到角两边距离相等：
 因 $AI$ 平分 $\angle BAC$，所以 $IE=IF$；
 因 $BI$ 平分 $\angle ABC$，所以 $IF=ID$。
 于是 $ID=IE=IF=r$。
 \item \textbf{计算三个小三角形。}由三角形面积等于底乘高的一半，
 \[
 \begin{aligned}
 S_{\triangle IBC}&=\frac12\,BC\cdot ID=\frac12ar,\\
 S_{\triangle ICA}&=\frac12\,CA\cdot IE=\frac12br,\\
 S_{\triangle IAB}&=\frac12\,AB\cdot IF=\frac12cr.
 \end{aligned}
 \]
 \item \textbf{合并并变形。}
 三块恰好组成 $\triangle ABC$，故
 \[
  S=\frac12ar+\frac12br+\frac12cr
   =\frac12r(a+b+c)=rp,
  \qquad \boxed{r=\frac Sp}.
 \]
 \item \textbf{顺便得到面积比。}
 三块面积都含有相同的系数 $r/2$，
 \[
   \boxed{S_{\triangle IBC}:S_{\triangle ICA}:S_{\triangle IAB}=a:b:c}.
 \]
 这说明求三小块面积比时连 $r$ 也不用求。
\end{enumerate}

\medskip
\par\noindent\begin{minipage}{\linewidth}
\textbf{\textcolor{CProof}{补充：反过来判断某点是不是内心}}

\PAngleFigure
\end{minipage}\par
图中把 $I$ 换成待判断的点 $P$：若两条垂线段相等，由共同斜边和一条直角边相等推出全等，再得两角相等。

若 $P$ \textbf{在三角形内部}，且它到 $AB,AC,BC$ 的垂直距离相等，
则 $P$ 是内心。理由是：从 $P$ 向 $AB,AC$ 作垂线，
两个直角三角形以 $AP$ 为共同斜边，另一条对应直角边又相等，
所以它们全等，推出 $AP$ 平分 $\angle BAC$；
对顶点 $B$ 同理，$BP$ 也平分 $\angle ABC$。
因此 $P$ 是两条内角平分线交点，即内心。
\textbf{特别强调：}“在内部”不能省去。

\medskip
\par\noindent\begin{minipage}{\linewidth}
\textbf{\textcolor{CProof}{切线段长的推导：三边拆成三对等长线段}}

\PTangentFigure{x}{y}{z}
\end{minipage}\par
先说明为什么同一顶点引出的两段相等：$IE\perp AE$、$IF\perp AF$，
两个直角三角形 $\triangle AEI$、$\triangle AFI$ 共用斜边 $AI$，
且 $IE=IF=r$，由斜边及一条直角边相等可知全等，故 $AE=AF=x$。
在 $B,C$ 处同理，$BF=BD=y$、$CD=CE=z$。

从图上沿每条边相加，得到
\[
 \underbrace{y+z}_{BC}=a,\qquad
 \underbrace{x+z}_{CA}=b,\qquad
 \underbrace{x+y}_{AB}=c.
\]
三式相加得 $2(x+y+z)=a+b+c=2p$，所以 $x+y+z=p$。
再分别减去 $y+z=a$、$x+z=b$、$x+y=c$，即得
\[
 \boxed{x=p-a},\qquad \boxed{y=p-b},\qquad \boxed{z=p-c}.
\]
因此三边已知时，\textbf{不必求半径或内心坐标}，就能确定切点把各边分成的长度。

\medskip
\textbf{\textcolor{CProof}{三种面积证明怎样选？}}

第一次学或者不理解等距性质时，用\textbf{方法一}；
考试中直接求面积、半径时，用\textbf{方法二}；
书面完整证明或需要同时说明面积比时，用\textbf{方法三}。
\end{proofbox}
